\documentclass[twocolumn]{notes}
\usepackage{notes}

\begin{document}
\title{Metrics to evaluate enhanced sampling}
\date{\today} % \formatdate{10}{11}{2022}
\maketitle
\tableofcontents
\newpage

See \texttt{plan.pdf} to see a good discussion of this.  Ultimately, just like the ACF is a ideal metric for CV quality, it is also a very good metric for implied timescale of the enhanced dynamics:
\begin{enumeratec}
    \item Separately from enhanced dynamics, find ideal CVs
    \item Now, use these CVs to assess \emph{apparent} timescale (from autocorrelation) of a segment of simulation data.  The quantity of such data will be very poor, so statistical noise will be high.  But it gives some indication of a true boost factor.
\end{enumeratec}


% \appendix
\end{document}



%% \href{https://www.ferglab.com/}{\color{darkred}Website}

% \usepackage{minted}
% \begin{minted}[mathescape, linenos]{python}
%     # Note: $\pi=\lim_{n\to\infty}\frac{P_n}{d}$
%     title = "Hello World"
% \end{minted}

%%=======================================================================
%%     PREFERRED FIGURE/TABLE
%%=======================================================================
%%                                 %%-------------------------------------%%
%% \begin{figure}[h!]              %%  caption right of figure
%%   \begin{minipage}[c]{0.67\textwidth}
%%     \includegraphics[width=\textwidth]{imgfile.png}
%%   \end{minipage}\hfill
%%   \begin{minipage}[c]{0.3\textwidth}
%%     \caption{Caption on the right hand side}
%%     \label{fig:}
%%   \end{minipage}
%% \end{figure}
%%                                 %%-------------------------------------%%
%% \begin{figure}[h!]\centering    %%  caption below figure
%%   \includegraphics[width=0.5\textwidth]{imgfile.png}
%%   \caption{General caption}
%%   \label{fig:label}
%% \end{figure}
%%                                 %%-------------------------------------%%
%% \begin{figure}[!h]\centering    %%  array of figures with subcaptions
%%   \subfloat[]{\includegraphics[width=0.5\textwidth]{}}
%%   \subfloat[]{\includegraphics[width=0.5\textwidth]{}} \\
%%   \subfloat[]{\includegraphics[width=0.5\textwidth]{}}
%%   \subfloat[]{\includegraphics[width=0.5\textwidth]{}} \\
%%   \subfloat[]{\includegraphics[width=0.5\textwidth]{}}
%%   \subfloat[]{\includegraphics[width=0.5\textwidth]{}}
%%   \caption{General caption}
%%   \label{fig:label}
%% \end{figure}
%%                                        %%------------------------------%%
%% \begin{table}[!h]\centering\ra{1.3}    %%  simple table
%%   \begin{tabular}{@{}rrrr@{}}\toprule
%%     & $t=0$ & $t=1$ & $t=2$ \\ \midrule
%%     $c$ & 1 & 2 & 3 \\
%%     $c$ & 1 & 2 & 3 \\
%%     \bottomrule
%%   \end{tabular}
%%   \caption{Caption}
%% \end{table}
%%                                        %%------------------------------%%
%% \begin{table}[!h]\centering\ra{1.3}    %%  complicated table
%%   \begin{tabular}{@{}rrrrcrrr@{}}\toprule
%%     & \multicolumn{3}{c}{$w = 8$} & \phantom{abc}&
%%     \multicolumn{3}{c}{$w = 16$}\\
%%     \cmidrule{2-4} \cmidrule{6-8} \cmidrule{10-12}
%%     & $t=0$ & $t=1$ & $t=2$ && $t=0$ & $t=1$ & $t=2$ \\ \midrule
%%     \rowcolor[gray]{.95} $c$ & 1 & 2 & 3 && 4 & 5 & 6 \\
%%     \rowcolor[gray]{1.0} $c$ & 1 & 2 & 3 && 4 & 5 & 6 \\
%%     \bottomrule
%%   \end{tabular}
%%   \caption{Caption}
%% \end{table}

%%=======================================================================
%%     OLDER FIGURE/TABLE
%%=======================================================================
%% \begin{center}
%%   \includegraphics[width=0.75\linewidth]{}
%%   \captionof{figure}{}
%% \end{center}
%%
%% \begin{center}
%%   \begin{tabular}{rl}
%%     \includegraphics[width=0.5\linewidth]{} &
%%     \includegraphics[width=0.5\linewidth]{} \\
%%     \includegraphics[width=0.5\linewidth]{} &
%%     \includegraphics[width=0.5\linewidth]{} \\
%%   \end{tabular}
%%   \captionof{figure}{}
%% \end{center}
%%
%% \begin{center}
%%   \begin{tabular}{lll}
%%     \hline
%%     A & B & C \\
%%     \hline
%%     \hline
%%     1 & 2 & 3 \\
%%     \hline
%%   \end{tabular}
%%   \captionof{table}{test}
%% \end{center}
